\begin{definition}[Multiplicative polynomial weights]\label{mod:finitefield-polynomialweights}
For a commutative semiring $R$, let
$\operatorname{MonicPolynomial}(R)\subset R[X]$ be the submonoid of monic
polynomials.  If $M$ is a monoid, define a monic-polynomial weight to be a
monoid homomorphism
\[
 \operatorname{MonicWeight}(R,M)
   =\operatorname{MonicPolynomial}(R)\longrightarrow M.
\]
A function $w:R[X]\to M$ with $w(1)=1$ which is multiplicative whenever
both inputs are monic induces such a weight.

Use Mathlib's $\operatorname{nextCoeff}(P)$ for the coefficient immediately
below the leading coefficient, and define
\[
 \operatorname{secondCoeff}(P)=[X^2]P^{\mathrm{reverse}}.
\]
Thus, if $2\leq\deg P$, then
\[
 \operatorname{secondCoeff}(P)
   =[X^{\deg P-2}]P,
\]
whereas $\operatorname{secondCoeff}(P)=0$ when $\deg P<2$.  In particular
this implements the manuscript's convention $a_2=0$ in degree one.
For monic $P,Q\in R[X]$, Mathlib's polynomial coefficient API gives
\[
 \begin{aligned}
  \operatorname{nextCoeff}(PQ)
    &=\operatorname{nextCoeff}(P)+\operatorname{nextCoeff}(Q),\\
  \operatorname{secondCoeff}(PQ)
    &=\operatorname{secondCoeff}(P)+\operatorname{secondCoeff}(Q)
      +\operatorname{nextCoeff}(P)\operatorname{nextCoeff}(Q).
 \end{aligned}
\]

Let now $k$ be a field, let
$\chi\in\operatorname{FiniteMulChar}(k)$, and let
$\psi\in\operatorname{FiniteAddChar}(k)$.  The Hasse--Davenport polynomial
weight is the element of $\operatorname{MonicWeight}(k,\mathbf C)$ given by
\[
 w_{\chi,\psi}(P)
 =\chi\!\left((-1)^{\deg P}[X^0]P\right)
    \psi\!\left(-\operatorname{nextCoeff}(P)\right).
\]
This is the manuscript's
$\chi((-1)^r a_r)\psi(-a_1)$, including its sign and constant-term
normalization.

Finally, suppose $\phi:k\to\mathbf C$ is nowhere zero and satisfies
\[
 \phi(x+y)=\phi(x)\phi(y)\psi(xy)
 \qquad(x,y\in k).
\]
These hypotheses imply $\phi(0)=1$.  The Hasse-function polynomial weight
is the element of $\operatorname{MonicWeight}(k,\mathbf C)$ given by
\[
 w_{\phi,\psi}(P)
 =\phi\!\left(-\operatorname{nextCoeff}(P)\right)
    \psi\!\left(-\operatorname{secondCoeff}(P)\right).
\]
Its multiplicativity uses the displayed second-coefficient identity: the
factor $\psi(\operatorname{nextCoeff}(P)\operatorname{nextCoeff}(Q))$
from the quadratic-refinement law cancels with the corresponding negative
factor from the second coefficient.
\LeanFile{LanglandsFirstMainLemma/FiniteField/PolynomialWeights.lean}
\lean{LanglandsFirstMainLemma.MonicPolynomial}
\lean{LanglandsFirstMainLemma.MonicWeight}
\lean{LanglandsFirstMainLemma.MonicWeight.ofFun}
\lean{LanglandsFirstMainLemma.secondCoeff}
\lean{LanglandsFirstMainLemma.monic_firstCoeff_mul}
\lean{LanglandsFirstMainLemma.monic_secondCoeff_mul}
\lean{LanglandsFirstMainLemma.gaussPolynomialWeight}
\lean{LanglandsFirstMainLemma.hasseFunction_zero}
\lean{LanglandsFirstMainLemma.hassePolynomialWeight}
\leanok
\SourceText{Proofs of Theorems thm:HD-lift and thm:Hasse-lift.}
\uses{mod:finitefield-characterapi}
\FormalizationContract
\end{definition}
